\begin{table}[!htbp]
\centering
\sbox{\tabletempbox}{%
\footnotesize
\begin{tabular}[t]{rrrr>{\raggedright\arraybackslash}p{15em}}
\toprule
Percentile & Actual $t(\hat{\alpha})$ & Simulated Mean & Prob.\ Luck & Interpretation\\
\midrule
1 & -3.619 & -2.064 & 0.0\% & Worse than luck (significant)\\
5 & -2.684 & -1.370 & 0.0\% & Worse than luck (significant)\\
10 & -2.184 & -1.053 & 0.0\% & Worse than luck (significant)\\
50 & -0.606 & -0.015 & 0.8\% & Indistinguishable from luck\\
90 & 0.993 & 1.014 & 50.1\% & Consistent with zero-skill\\
95 & 1.362 & 1.326 & 59.6\% & Consistent with zero-skill\\
99 & 2.248 & 2.019 & 80.9\% & Consistent with zero-skill\\
\bottomrule
\end{tabular}%
}
\setlength{\tabletempwidth}{\wd\tabletempbox}
\ifdim\tabletempwidth>\linewidth\setlength{\tabletempwidth}{\linewidth}\fi
\begin{minipage}{\tabletempwidth}
\captionsetup{width=\linewidth}
\caption{\label{tab:bootstrap_tails}Fama--French (2010) Bootstrap: Actual vs.\ Simulated $t(\hat{\alpha})$ Percentiles}
\ifdim\wd\tabletempbox>\linewidth
  \resizebox{\linewidth}{!}{\usebox{\tabletempbox}}
\else
  \usebox{\tabletempbox}
\fi
\par\medskip
\begin{singlespace}\footnotesize\noindent
Bootstrap procedure follows \textcite{FamaFrench2010}. Sample: actively managed funds, Incubation-corrected (Evans 2010) panel (no date cap), performance-comparison subsample per flagged\_funds.xlsx, minimum 24 monthly observations ($N = 1815$ funds). For each fund, estimated monthly alpha is subtracted from the excess return series to construct a zero-alpha null return. In each of $B = 10{,}000$ bootstrap iterations, calendar months are resampled with replacement, preserving cross-sectional factor return dependence. The \textcite{Carhart1997} four-factor model is re-estimated on each resampled series. \textit{Actual} $t(\hat{\alpha})$: percentile of the empirical $t$-statistic distribution across active funds. \textit{Simulated Mean}: average of that percentile across all iterations. \textit{Prob.\ Luck}: fraction of iterations in which the simulated percentile falls below the actual value; values below 5\% at lower percentiles indicate underperformance unlikely to be explained by luck alone. Newey-West standard errors with a 6-month lag are used throughout.
\end{singlespace}
\end{minipage}
\end{table}

